\documentclass[11pt]{article}
\usepackage[a4paper,margin=22mm]{geometry}
\usepackage{amsmath,amssymb,amsthm,mathtools,bm}
\usepackage{booktabs,array,microtype,xurl,hyperref}
\hypersetup{colorlinks=true,linkcolor=blue,citecolor=blue,urlcolor=blue,
pdftitle={P54 complex phases and Pati--Salam matching}}
\setlength{\emergencystretch}{2em}
\setcounter{tocdepth}{1}
\newcommand{\tr}{\operatorname{Tr}}
\newcommand{\MS}{\overline{\mathrm{MS}}}
\newcommand{\diag}{\operatorname{diag}}
\newcommand{\norm}[1]{\left\lVert#1\right\rVert_2}
\newtheorem{proposition}{Proposition}
\title{Route-F P54: Common Complex Phases,\\Finite Pati--Salam Matching,\\and the Self-Consistency Interface}
\author{P54PQ-v2 four-dimensional mainline}
\date{5 September 2026}
\begin{document}
\maketitle
\begin{abstract}
We transport the actual four-copy scalar eigenvector into a single
standard-hypercharge spinor convention. This fixes the up/down conjugation,
the relative Majorana/triplet phases, and the interference of the two seesaw
terms. The previous magnitude-only light-copy dictionary is superseded.
We also construct a nearby Pati--Salam-symmetric stationary saddle of the
same frozen action, whose integrated 55-real-dimensional scalar block is
positive. It yields a finite upper gauge threshold, including the vector
constant. Retained tachyons do not invalidate this heavy-field matching.
The lower effective-action matching and finite Yukawa diagrams must still
be supplied before a physical global fit is defined. Numerical algebraic
regressions and synthetic family choices are not experimental fits.
\end{abstract}
\tableofcontents
\clearpage

\section{Scope and the fixed action}
The action, field content and scalar benchmark remain P54PQ-v2:
$3\,16_F+54_{H,\mathbb R}+126_{H,\mathbb C}+10_{H,\mathbb C}+1_{H,\mathbb C}$,
with PQ charges $(-1,0,+2,-2,-4)$ in the order
$(\Psi,\Phi,\Sigma,\phi,S)$. No extra family matrix, scalar multiplet,
torsion sector or lattice scan is introduced. The previous checkpoint was
committed and pushed as \texttt{b8e7398}; the calculations below are the
subsequent checkpoint.

We retain 328 canonical real scalar coordinates and the fixed-VEV
background-field Landau-$\MS$ curvature prescription from the companion
\texttt{p54\_fixed\_vev\_full\_doublet\_matching.tex}. Numerical masses below
are in units of the historical $\omega=1$. The old physical values of
$M_I,M_U$ remain unrecertified: the earlier threshold census was inconsistent
with its PS beta functions. New dimensionless matching tensors do not by
themselves supply replacement physical scales.

\section{One complex scalar--spinor dictionary}
\subsection{Construction from Clifford intertwiners}
On five qubits define
\begin{align}
 \Gamma_{2j-1}&=Z^{\otimes(j-1)}\otimes X\otimes I^{\otimes(5-j)},&
 \Gamma_{2j}&=Z^{\otimes(j-1)}\otimes Y\otimes I^{\otimes(5-j)},\\
 \Gamma_*&=(-i)^5\Gamma_1\cdots\Gamma_{10},&
 C&=\Gamma_2\Gamma_4\Gamma_6\Gamma_8\Gamma_{10}.
\end{align}
They obey $\{\Gamma_i,\Gamma_j\}=2\delta_{ij}$ and $C^T=-C$.
The restrictions of $C\Gamma_i$ and $C\Gamma_{[5]}$ to the appropriate
chiral subspace are symmetric, as required by two left-handed Weyl fields.
For an antisymmetric vector generator $R$,
\begin{equation}
 S(R)=\frac12\sum_{i<j}R_{ij}\Gamma_i\Gamma_j,
 \qquad S(R)^T B_a+B_a S(R)+\sum_bB_b R_{ba}=0.
\end{equation}
The five-form version uses the induced wedge representation. The numerical
audit tests all 45 generators, not only the hypercharge eigenvalues.

Let $U_{126}$ be the existing orthonormal ordered-five-form embedding.
Its old orientation selects the conjugate standard-matter hypercharges.
We make an explicit global coordinate change, rather than conjugating an
isolated Clebsch coefficient:
\begin{equation}
 x_{\rm phys}=Kx_{\rm old},\qquad U_{126,\rm phys}=\overline U_{126},
 \qquad \Gamma_*\Psi_{\rm phys}=+\Psi_{\rm phys}.
\end{equation}
Here $K$ flips the imaginary real-coordinate blocks of $\Sigma,\phi,S$ and
leaves all other coordinates unchanged. Consequently
\begin{equation}
 V_{\rm phys}(x)=V_{\rm old}(Kx),\quad H_{\rm phys}=KH_{\rm old}K,
 \quad B_{\rm phys}=KB_{\rm old}.
\end{equation}
If complex old-action couplings had been fixed, their conjugation is
transported too. In the new raw convention
\begin{equation}
 -\mathcal L_Y=\frac12h_{ij}\Psi_i^TC\Gamma_a\Psi_j\phi_a^*
 +\frac1{2\cdot5!}f_{ij}\Psi_i^TC\Gamma_{a_1\cdots a_5}\Psi_j
 \Sigma_{a_1\cdots a_5}+\mathrm{h.c.}
\end{equation}
The physical spinor has hypercharge multiset
$(1/6)^6,(-2/3)^3,(1/3)^3,(-1/2)^2,1,0$.

The actual old copy order is retained:
\begin{equation}
 B=(B_{\phi,\rm hol},B_{\phi,\rm anti},
 B_{\Sigma,\rm hol},B_{\Sigma,\rm anti}),\qquad
 x_{\rm phys}=2\operatorname{Re}(B_{\rm phys}c H^0),\quad
 \langle H^0\rangle=v/\sqrt2.
\end{equation}
All right-handed-conjugate component phases are fixed once by requiring
the corresponding reference $10_H$ coefficient to be positive. The same
states are then used for every $126_H$, Majorana and triplet entry. Direct
matrix-element contraction gives
\begin{equation}
\begin{array}{c|cc}
 &h_{\rm raw}&f_{\rm raw}\\\hline
 Y_u&\sqrt2\,c_1&2c_4/\sqrt3\\
 Y_d&\sqrt2\,c_2^*&2c_3^*/\sqrt3\\
 Y_e&\sqrt2\,c_2^*&-6c_3^*/\sqrt3\\
 Y_\nu&\sqrt2\,c_1&-6c_4/\sqrt3
\end{array}
\label{eq:copy}
\end{equation}
These are matrix coefficients: each entry multiplies its raw family matrix.
The $c$ versus $c^*$ distinction is required by $H^0$ versus $H^{0*}$.
It cannot be inferred from $|c_a|$.

\subsection{The corrected bosonic light map}
Define Dirac-unit couplings $h_D=\sqrt2h_{\rm raw}$ and
$f_D=2f_{\rm raw}/\sqrt3$. Equation~\eqref{eq:copy} implies
\begin{align}
Y_d&=a h_D+d f_D,&Y_e&=a h_D-3d f_D,\\
Y_u&=b h_D+e f_D,&Y_\nu&=b h_D-3e f_D,\qquad
(a,b,d,e)=(c_2^*,c_1,c_3^*,c_4).
\end{align}
At the fixed-VEV bosonic-retuned point,
\begin{align}
a&=-0.7631652014906-1.03\times10^{-12}i,&b&=0.6462011667549,\\
d&=1.2183583687513\times10^{-5}+5.47\times10^{-15}i,&
e&=-0.001710897393866-1.02\times10^{-14}i.
\end{align}
For nonzero $a,b,d$ the conventional ratios are
\begin{equation}
r=b/a=-0.846738249455,\qquad s=ae/(bd)=165.843989136.
\end{equation}
The same dictionary at the tree point gives
$r=-0.905014358012$, $s=-29.7554251165$.
In particular the small coefficient $d$ crosses zero under the bosonic
correction. A large change in $s$, which divides by $d$, is not a
large rotation of the entire scalar eigenvector. Numerically stable fits
must use $(a,b,d,e)$ directly, not $(r,s)$ near $d=0$.
Earlier reciprocal/magnitude-only copy maps, and bounds using those maps,
are not current physical constraints. This is not a claim that flavor
tension has disappeared, nor a whole-model exclusion.

\subsection{Majorana and triplet phases in the same basis}
The canonical singlet and LL-triplet contractions are
\begin{equation}
 \kappa_R=+i4\sqrt2,\qquad \kappa_{LL}=-i4\sqrt2,\qquad
 f_M=\kappa_R f_{\rm raw}=i2\sqrt6 f_D,
 \qquad M_R=\sigma\omega f_M.
\end{equation}
Here the numerical $\sigma=0.1265$ is dimensionless and $\omega$ restores
mass units. The $i$ is a component-basis convention, not an additional
physical CP parameter. Dropping it in only one term would nevertheless
change the declared matching map.

The existing two-triplet tree Hessian and cubic source, evaluated along the
bosonic-improved light direction, induce
\begin{equation}
 \frac{\omega\Delta_{126}}{v^2}=-0.003323921052868.
\end{equation}
The $54$ triplet has zero renormalizable LL coupling but contributes to this
response through its mixing with $126$. Contracting the full response in
the common dictionary yields
\begin{align}
 M_\nu&=M_L-\frac{v^2}{2}Y_\nu M_R^{-1}Y_\nu^T\\
 &=\frac{v^2}{\omega}
 \left[C_{II,D}f_D+C_{I,D}Y_\nu f_D^{-1}Y_\nu^T\right],\\
 C_{II,D}&=0.01628382104964i,&
 C_{I,D}&=\frac{-1}{2\sigma i2\sqrt6}=0.80681480328827i.
\end{align}
This formula requires invertible $f_D$. If a neutrino is not hard at the
matching scale it must remain in the EFT; the inverse is not to be
regularized by hand. The scalar response is still mixed-order
(tree $K,J$ with bosonic one-loop $c$), not complete one-loop Weinberg
matching.

For independent neutrino basis changes $\nu_L=U_L\nu_L'$,
$\nu^c=U_R\nu^{c\prime}$, one has
\begin{equation}
Y_\nu'=U_L^TY_\nu U_R,\quad M_R'=U_R^TM_RU_R,\quad
M_L'=U_L^TM_LU_L,
\end{equation}
so the full seesaw transforms as $M_\nu'=U_L^TM_\nu U_L$.
The proof follows by inserting
$(U_R^TM_RU_R)^{-1}=U_R^\dagger M_R^{-1}U_R^*$.
The verifier also rotates actual scalar copies by random $U(4)$ and tests
the conjugate down rule. Deliberately replacing down $c^*$ with $c$ gives
an order-one residual, so the regression is not insensitive to this error.

\section{A controlled upper PS finite gauge threshold}
\subsection{Same-action stationary saddle and heavy-field elimination}
Retain the four complex PS parents
\begin{equation}
 (1,2,2)_{10},\quad(15,2,2)_{126},\quad
 (\overline{10},3,1)_{126},\quad(10,1,3)_{126}.
\end{equation}
These are the common standard-matter labels, with
$4=(3,1/3)+(1,-1)$ under $SU(3)\times U(1)_{B-L}$.
The old literal-P1 labels are conjugated by the dictionary change.
The LL scalar has $B-L=+2$ and its singlet triplet has $Y=+1$;
the right-handed-conjugate pair instead requires scalar $B-L=-2$.
Their Casimirs and real index traces are unchanged.
Their real dimension is $8+120+60+60=248$.
At $\sigma=0$ solve the two heavy radial equations
\begin{equation}
 \partial_\omega V(\omega,0,v_s)=0,\qquad
 \partial_{v_s}V(\omega,0,v_s)=0
\end{equation}
with every existing parameter, including the three quadratic masses,
unchanged. This gives
\begin{equation}
 \omega_{PS}=1.0008103059,\qquad v_{s,PS}=0.2542776971.
\end{equation}
The full gradient vanishes. There are 24 upper gauge Goldstones and one PQ
direction. The complement of these and the active 248-plane has dimension
55 and satisfies
\begin{equation}
\lambda_{\min}(H_{HH})=0.10288204>0.
\end{equation}
All 124 negative directions are retained active fields, rather than modes
being integrated out. Thus this point is a stationary saddle, not a proposed
stable vacuum of the full theory.

\begin{proposition}[Local heavy-field EFT]
On a fixed gauge slice, let $h$ be heavy and $\ell$ retained coordinates.
If $V_h(h_0,\ell_0)=0$ and $V_{hh}>0$, the implicit-function theorem yields
a unique local heavy valley $h=h_*(\ell)$ and a real local tree EFT
$V_{\rm EFT}(\ell)=V(h_*(\ell),\ell)$. No positivity assumption on the
retained Hessian is needed.
\end{proposition}
\begin{proof}
Invertibility of $V_{hh}$ applies the implicit-function theorem to $V_h=0$.
Differentiating gives $h_*'=-V_{hh}^{-1}V_{h\ell}$. At a stationary point,
\begin{equation}
 (V_{\rm EFT})_{\ell\ell}
 =V_{\ell\ell}-V_{\ell h}V_{hh}^{-1}V_{h\ell}.
\label{eq:schurEFT}
\end{equation}
The positive heavy block fixes the local branch and its real Gaussian
functional. Tachyons in $\ell$ are represented by the EFT itself.
\end{proof}
This is local existence, not a global heavy valley or a proof that a
renormalizable truncation captures all powers of $\sigma/\omega$.

\subsection{Finite constants and logarithmic consistency}
Use real-representation index operators normalized to
$\tr I_i=T_i(R_{\mathbb R})$. For the present upper threshold, with no
massive charged fermions removed, the gauge matching is \cite{Hall}
\begin{align}
 \alpha_i^{-1}(\mu)&=\alpha_{10}^{-1}(\mu)-\frac{\lambda_{U,i}(\mu)}{12\pi},\\
 \lambda_{U,i}&=\tr_{V_H} I_i
 -21\tr_{V_H}\!\left[I_i\log\frac{M_V}{\mu}\right]
 +\tr_{S_H}\!\left[I_i\log\frac{M_S}{\mu}\right].
\label{eq:finiteGauge}
\end{align}
Goldstone directions are removed once from the physical scalar trace; the
massive vector coefficient includes the corresponding gauge/ghost sector.
The finite vector constant cannot be recovered by matching logarithmic
slopes alone. At the PS saddle the active, mass and Goldstone projectors
commute. Thus the heavy traces are genuine spectral restrictions.

At $\mu_U=g_U\omega_{PS}$, in group order $(4,L,R)$,
\begin{equation}
 \lambda_U=(2.944052689,\ 7.602978830,\ 7.602978830),
 \qquad \lambda_U^{V,\rm const}=(4,6,6).
\end{equation}
The upper physical real scalar indices are $(12,6,6)$, so
\begin{equation}
 \lambda_U'=21(4,6,6)-(12,6,6)=(72,120,120).
\end{equation}
These equal $-6(a_{10}-a_{PS})$ with
$a_{10}=-34/3$ and $a_{PS}=(2/3,26/3,26/3)$, as follows directly by
differentiating the inverse-coupling matching equation.
The largest retained absolute squared mass divided by the smallest
integrated squared mass is $1.01260824$: there is no uniform hierarchy
between every retained and integrated scalar. Positivity defines the local
heavy functional, but does not prove a controlled two-site renormalizable
truncation for all intermediate loop momenta.

At the lower site the repaired census fixes, in SM order $(3,2,1)$,
\begin{equation}
 I_V=(1,0,14/5),\quad I_S=(67,71,67),\quad
 \lambda_I'=(-46,-71,-41/5).
\end{equation}
The hypercharge projection is
$a_1^{PS}=2a_4/5+3a_R/5$. These three lower slopes, together with the three
upper slopes, reproduce all six beta discontinuities. This does not
calculate the finite lower constants.

At the simultaneous broken vacuum the old active projector $P_A$ obeys
\begin{equation}
 \lVert[P_A,H]\rVert_F\simeq0.265322,\qquad
 \lVert[P_A,P_{G_U}]\rVert_F\simeq0.609960,
 \qquad\tr(P_A P_{G_U})\simeq0.189003.
\end{equation}
These are structural violations, not eigensolver rounding. A trace of
$P_A\log H$ there cannot replace the lower EFT determinant.
As a separate check, the complete one-step $SO(10)\to SM$ hard spectrum gives
\begin{equation}
 \lambda_{\rm one\ step}=(36.63010816,-6.54752280,119.73803708),
 \quad\lambda'=(26,49,92.6).
\end{equation}
It is a full-spectrum regression, not an adopted change from the two-stage
physical branch. Finite shifts of tree masses from the one-loop tadpole
counterterms inside a one-loop threshold enter at higher order; all such
order choices and the background prescription must remain common when the
full finite EFT is assembled.

\subsection{The actual lower heavy valley is gapped, but is not a mass-log shortcut}
Transport the upper heavy 55-plane to the actual broken vacuum. Let its
orthonormal basis be $B_H$, let $B_A$ span the retained parent plane, and
remove the upper orbit from the latter:
\begin{equation}
B=(I-P_{G_U})B_A,\quad C=B_H^THB_H,\quad
A=B^THB,\quad W=B_H^THB,\quad E=-C^{-1}W.
\end{equation}
The calculation verifies $B_H^TB=0$, SM covariance, and
$\lambda_{\min}(C)=0.099450257411$. The tree heavy-valley action has
\begin{equation}
S=A-W^TC^{-1}W,\qquad G=B^TB+E^TE.
\end{equation}
The generalized problem $Sv=m^2Gv$ has no negative eigenvalues and exactly
13 zero directions (nine lower Goldstones and four real Higgs components).
Its Ward residual is $3.21\times10^{-16}$, and
\begin{equation}
\lambda(G)\subset[0.984249788817,1.01800727050],\qquad
\norm{E}=0.134191171461.
\end{equation}
This verifies a local derivative-expansion interface on the actual vacuum.
It does not make these generalized eigenvalues full pole masses. Indeed,
the exact Euclidean quadratic Schur operator is
\begin{align}
\mathcal K_{\rm EFT}(p^2)
 &=p^2G_0+A-W^T(p^2I+C)^{-1}W,\quad G_0=B^TB,\\
 &=S+p^2G-p^4W^TC^{-3}W+O(p^6).
\end{align}
At the operator level
$\det\mathcal K_{\rm full}=\det\mathcal K_{HH}\det\mathcal K_{\rm EFT}$.
Gauging this identity also generates covariant vertices and commutators;
replacing it by $\tr\log S$ loses those terms. The implemented
two-state, equal-gauge-charge counterexample distinguishes the exact Schur
determinant from a projector-weighted mass logarithm. Completing the lower
finite matching means using the induced action, not merely another spectral
assignment rule.

\section{Complete four-parent one-loop PS Yukawa system}
\subsection{Derivation and normalization from the actual tensors}
Write the real-scalar Weyl Lagrangian as
$-\mathcal L_Y=\frac12\mathcal Y^a_{IJ}\psi_I\psi_J q_a+\mathrm{h.c.}$,
where each full $\mathcal Y^a$ is symmetric. For the two-component
convention, the general one-loop coefficient is \cite{LWX}
\begin{align}
16\pi^2\beta_{\mathcal Y^a}
={}&\tfrac12(S^T\mathcal Y^a+\mathcal Y^aS)
 +2\sum_b\mathcal Y^b\mathcal Y^{a\dagger}\mathcal Y^b\nonumber\\
&+\sum_b\mathcal Y^b\operatorname{Re}\tr(\mathcal Y^{b\dagger}\mathcal Y^a)
 -3\sum_i g_i^2(C_i^T\mathcal Y^a+\mathcal Y^a C_i),\qquad
 S=\sum_b\mathcal Y^{b\dagger}\mathcal Y^b.
\end{align}
Every complex scalar supplies a pair
$\mathcal Y^{\rm Re}=\mathcal Y/\sqrt2$ and
$\mathcal Y^{\rm Im}=\pm i\mathcal Y/\sqrt2$. For each internal pair,
$\mathcal Y^{\rm Re}\mathcal Y^{a\dagger}\mathcal Y^{\rm Re}
+\mathcal Y^{\rm Im}\mathcal Y^{a\dagger}\mathcal Y^{\rm Im}=0$.
This cancels the vertex sum without deleting any scalar parent.

The actual 124 complex PS intertwiners are projections of the common
Clifford tensors, not arbitrary Clebsch tables. In a PS fermion order
$(\psi_L,\psi_R)$, a bidoublet family matrix $H$ enters opposite off-diagonal
blocks as $H$ and $H^T$; thus the full Weyl tensor is symmetric even when
$H$ itself is not. The canonical PS normalizations at the tree upper
boundary are
\begin{equation}
H=\sqrt2h_{\rm raw},\quad F=4f_{\rm raw},\quad
L=R=2\sqrt2f_{\rm raw}.
\label{eq:PSboundary}
\end{equation}
In particular $F=2\sqrt3 f_D$ and $f_M=i2R$ in the common phase convention.
Different normalizations in a reference must be converted before importing
coefficients. The $\Delta_R$-only system in Ref.~\cite{MOR} is a comparison
limit, not this active census.

Define
\begin{align}
A_L&=HH^\dagger+\tfrac{15}{4}FF^\dagger+\tfrac{15}{2}LL^\dagger,\\
A_R&=H^\dagger H+\tfrac{15}{4}F^\dagger F+\tfrac{15}{2}R^\dagger R,\\
G_D&=\tfrac94(g_L^2+g_R^2+5g_4^2),&
G_L&=\tfrac94(2g_L^2+5g_4^2),&G_R&=\tfrac94(2g_R^2+5g_4^2).
\end{align}
The complete matrix equations are
\begin{align}
16\pi^2\dot H&=A_LH+HA_R+[4\tr(H^\dagger H)-G_D]H,\\
16\pi^2\dot F&=A_LF+FA_R+[\tr(F^\dagger F)-G_D]F,\\
16\pi^2\dot L&=A_LL+LA_L^T+[2\tr(L^\dagger L)-G_L]L,\\
16\pi^2\dot R&=A_R^TR+RA_R+[2\tr(R^\dagger R)-G_R]R,
\qquad\dot{}=d/d\log\mu.
\label{eq:PSbeta}
\end{align}
Here $L=L^T,R=R^T$ are preserved. Away from the appropriate parity locus,
$H,F$ are general complex matrices; projecting their antisymmetric parts
away at every step would change the differential equation.
The actual 248-real-scalar, 48-Weyl-component tensors reproduce
Eq.~\eqref{eq:PSbeta} component by component for unequal left/right inputs.

The associated Yukawa contribution to the two-loop gauge flow is
\begin{equation}
Y_4=\big(2\tr(A_L+A_R),\ 4\tr A_L,\ 4\tr A_R\big),
\quad
\dot g_i=\frac{g_i^3a_i}{16\pi^2}
+\frac{g_i^3}{(16\pi^2)^2}\left[\sum_jb_{ij}g_j^2-Y_{4,i}\right],
\end{equation}
with
\begin{equation}
a=(2/3,26/3,26/3),\qquad
b=\begin{pmatrix}3551/6&249/2&249/2\\
1245/2&779/3&48\\1245/2&48&779/3\end{pmatrix}.
\end{equation}
For example the SU(4) factor follows from $C_4(4)=15/8$, eight left and
eight right gauge components, and $\dim SU(4)=15$ in
$\tr(C_iS)/\dim G_i$. The implementation checks this trace using the
actual tensors and also includes a Standard-Model top-only normalization
regression.

\subsection{A protected one-loop submanifold, not new independent matrices}
There is a useful exact simplification of the frozen tree boundary.
Suppose
\begin{equation}
g_L=g_R,\qquad H=H^T,\quad F=F^T,\quad
L=R=F/\sqrt2.
\label{eq:protected}
\end{equation}
Then $A_R=A_L^T$, $G_D=G_L=G_R$, and
$\tr F^\dagger F=2\tr R^\dagger R$.
Inserting these identities into Eq.~\eqref{eq:PSbeta} gives
\begin{equation}
\beta_H=\beta_H^T,\qquad \beta_F=\beta_F^T,\qquad
\beta_F=\sqrt2\beta_L=\sqrt2\beta_R.
\end{equation}
The gauge beta functions also give $\beta_{g_L}=\beta_{g_R}$, since their
rows and $\tr A_L,\tr A_R$ agree on this locus. Uniqueness of the ODE
therefore proves invariance of \eqref{eq:protected}. This is a conditional
one-loop protection, not merely a numerical coincidence.

It follows that the tree Spin(10) boundary does \emph{not} acquire three
independent arbitrary $126$ family matrices just by PS running. On this
locus the ratio $f_M/f_D=i2\sqrt6$ is preserved at one loop. For finite
matching that moves the boundary off the locus, unequal left/right
thresholds, decoupling, or higher-loop evolution, the general system must
instead transport $F,L,R$ separately. The lower dictionary is always
\begin{equation}
Y_d=aH+\frac{d}{2\sqrt3}F,\quad
Y_u=bH+\frac{e}{2\sqrt3}F,\quad
M_R=i2\sigma\omega R,\quad M_L=-i2\Delta_L L,
\end{equation}
with the same lepton $-3$ replacements. All four PS matrices remain
correlated outputs of two UV matrices plus calculated thresholds. They
are not four newly fitted free matrices.

\section{Finite complex kinetic matching: an executable interface}
\subsection{Canonical normalization and the scalar curvature}
In the Weyl convention $\psi_L^T Y_a\psi_R z_a$, write the hard local
kinetics as $Z_L=I+K_L$, $Z_R=I+K_R$, $Z_H=I+K_H$ and the vertex as
$\widetilde Y_a=Y_a+\delta Y_a$. If these Hermitian metrics are positive,
set $S_j=Z_j^{-1/2}$. Direct field substitution gives
\begin{align}
 Y^{\rm can}_a&=\sum_b S_L^T\widetilde Y_b S_R(S_H)_{ba},\\
 D_{\rm can}&=S_H^\dagger D S_H,\\
 Y^{\rm can}_a&=Y_a+\delta Y_a-\frac12\left[
 K_L^TY_a+Y_aK_R+\sum_bY_b(K_H)_{ba}\right]+O(\hbar^2).
\label{eq:finiteY}
\end{align}
The linear formula is the general canonical matching algebra; see also
Appendix A of Ref.~\cite{PatelShukla}. Its model-specific finite matrices
must still be computed from UV-minus-EFT amplitudes. SU(5) group factors
from that paper are not P54 coefficients.

If $Dc=0$ and $c^\dagger c=1$, then
\begin{equation}
 c_{\rm can}=\frac{Z_H^{1/2}c}{\sqrt{c^\dagger Z_H c}},\qquad
 D_{\rm can}c_{\rm can}=0,\qquad
 S_Hc_{\rm can}=\frac{c}{\sqrt{c^\dagger Z_Hc}}.
\label{eq:kineticzero}
\end{equation}
Thus at an exact zero the original-coordinate light direction changes only
by kinetic normalization. An apparent rotation in canonical coordinates
is not an additional independent scalar-mixing correction. Vertices and
curvature must be transformed together to see this cancellation.
For down/e vertices multiplying $z^*$, use $Z_H^*$ and $c^*$, not the up
formula. The executable tests use the actual four-copy embedding and
curvature with synthetic complex finite matrices. They verify independent
left/right family transformations, arbitrary $U(4)$ copy rotations, and
the conjugate-copy rule.

\subsection{Stable finite kernels and what they do not calculate}
For squared masses $x,y\geq0$, at least one nonzero, define
\begin{align}
 f(x,y;\mu)&=-\frac1{16\pi^2}\int_0^1
 \log\frac{(1-t)x+ty}{\mu^2}\,dt,\\
 h(x,y;\mu)&=\frac1{16\pi^2}\int_0^1(1-t)
 \log\frac{(1-t)x+ty}{\mu^2}\,dt.
\end{align}
Integrating gives, for unequal positive masses,
\begin{align}
16\pi^2 f&=1-\frac{x\log(x/\mu^2)-y\log(y/\mu^2)}{x-y},\\
16\pi^2 h&=\frac12\log(x/\mu^2)
 +\frac{\frac12r^2\log r-\frac34r^2+r-\frac14}{(1-r)^2},\quad r=y/x.
\end{align}
The integral form avoids subtractive cancellation near equality. At
$x=y$, $(16\pi^2f,16\pi^2h)=(-\log(x/\mu^2),\frac12\log(x/\mu^2))$.
At $y=0$ they become
$(1-\log(x/\mu^2),\frac12\log(x/\mu^2)-\frac14)$; at $x=0$ the last
constant is $-3/4$. Also
\begin{equation}
 \frac{d(f,h)}{d\log\mu}=\frac{(2,-1)}{16\pi^2}.
\end{equation}
These kernels agree with Appendix B of Ref.~\cite{PatelShukla} and are
implemented with equal/unequal/one-zero mass regressions. They are building
blocks, not the complete tensor contractions of P54 diagrams. The API
rejects missing finite vertices, nonpositive kinetic metrics and purely
light kernels rather than filling absent physical inputs with zeros.
Using exact square roots of one-loop input matrices is an algebraic
regression, not a calculation of two-loop terms.

\section{Local scalar--flavor feedback and its analytic derivative}
\subsection{An actual Clebsch-constrained inner solve}
Two deterministic small complex symmetric raw family choices test the
inner scalar calculation. They are not measured or globally fitted
matrices. The actual copy dictionary supplies all $Y_{\nu,a}$ and
$M_R=\kappa_R\sigma f_{\rm raw}$ in $\omega=1$ units. With
$X=M_R^\dagger M_R$ and $W=X[\log(X/\mu^2)-1]$, use
\begin{align}
\Pi_{F,ab}&=-\frac1{8\pi^2}\tr(Y_{\nu,a}^\dagger Y_{\nu,b}W),\\
\delta\nu_F^2&=-\frac1{8\pi^2\sigma^2}\tr\{X^2[\log(X/\mu^2)-1]\},\\
\Delta_F&=\Pi_F-\tfrac12\delta\nu_F^2P_{126},\\
D(\theta,\xi)&=D_0+\Delta_B+\Delta_F(\theta)-\xi P_{10}.
\end{align}
The subtraction scale is held fixed in all field derivatives. The scalar
$\xi$ is tuned to the first stable zero of the complete complex matrix,
not just to a scalar projection. At a gapped zero $Dc=0$, $c^\dagger c=1$,
choose $c^\dagger c'=0$. Differentiation gives
\begin{equation}
\xi'=\frac{c^\dagger\Delta_F'c}{c^\dagger P_{10}c},\qquad
c'=-D^+(\Delta_F'-\xi'P_{10})c.
\label{eq:implicit}
\end{equation}
Here $D^+$ is the inverse on the heavy complement, zero on the light
direction. The formula is valid while the heavy gap and
$c^\dagger P_{10}c$ are nonzero.

For the explicit test $f_{\rm raw}(\theta)=(1+\theta)f_{\rm raw}$,
\begin{equation}
X'=2X,\quad W'=2X\log(X/\mu^2),\qquad
(\delta\nu_F^2)'=-\frac{\tr\{X^2[4\log(X/\mu^2)-2]\}}
{8\pi^2\sigma^2}.
\end{equation}
Apply the product rule also to both $Y_{\nu,a}$ in $\Pi_F$ to obtain
$\Delta_F'$ without a finite difference. Then Eq.~\eqref{eq:implicit}
agrees with two independent finite-difference steps of the retuned
eigenpair. Representative results are
\begin{center}
\begin{tabular}{@{}rrrr@{}}
\toprule
case&$\xi$&heavy gap/$\omega^2$&$\norm{dc/d\theta}$\\\midrule
0&$-0.02025044324$&$0.03026151664$&$0.0008092458$\\
1&$-0.02024289530$&$0.03026452622$&$0.0012951028$\\\bottomrule
\end{tabular}
\end{center}
The small-step relative error of $c'$ is about $6\times10^{-9}$.
This validates a differentiable local inner solver. It neither selects
physical family parameters nor replaces the finite PS map.

\subsection{Recompute the type-II source when the light direction changes}
The previous type-II number cannot be kept fixed while $c$ changes.
Using the quartic action, reconstruct the exact cubic jet tensor
\begin{equation}
T_{Apq}=V'''[e_A,q_p,q_q],\qquad
J_A(c)=\sum_{pq}T_{Apq}c_{\mathbb R,p}c_{\mathbb R,q},\quad
c_{\mathbb R}=(\operatorname{Re}c,\operatorname{Im}c).
\end{equation}
Centered differences of the Hessian are algebraically exact cubic jets
for a quartic potential. Eight already cached Hessians give the four real
Higgs-direction jets; hypercharge covariance supplies the four imaginary
ones. This introduces no new parameter or lattice scan and reproduces
the old source to $4.4\times10^{-16}$ relative error.
For each new eigenvector recompute
\begin{equation}
z_{\mathbb R}(c)=-\tfrac12K_{\rm tree}^{-1}J(c),\quad
\Delta_A(c)=\frac{z_{R,A}+iz_{I,A}}{\sqrt2},\quad
C_{II,\rm raw}(c)=\sum_A\kappa_{LL,A}\Delta_A(c).
\end{equation}
For the two examples $|\Delta C_{II,\rm raw}|$ relative to the old
bosonic direction is $2.295\times10^{-5}$ and $3.623\times10^{-5}$.
All inputs, including the cubic tensor, triplet Hessian, source and LL
coefficient, are transported in the CP-conjugation test. This tests
covariance, not CP invariance of a chosen complex benchmark.
The triplet Hessian and cubic vertices are still tree quantities;
their complete loop correction, evolved $L/R$ inputs off the protected
locus, and finite Weinberg matching remain necessary for a physical fit.

\section{Research interpretation and completion boundary}
The executable checkpoint passes the following tests on a fresh replay:
\begin{center}
\begin{tabular}{@{}lr@{}}
\toprule
Verifier suffix (all under \texttt{route\_f/code/})&passed\\\midrule
\texttt{common\_yukawa\_phase}&28/28\\
\texttt{ps\_finite\_thresholds}&25/25\\
\texttt{ps\_yukawa\_flow}&29/29\\
\texttt{finite\_yukawa\_interface}&17/17\\
\texttt{self\_consistent\_light}&36/36\\\bottomrule
\end{tabular}
\end{center}
The full filenames begin with \texttt{verify\_p54\_} and end in
\texttt{.py}; each exports a same-named JSON/Markdown ledger, with input
hashes and individual residuals, under \texttt{route\_f/output/}.
The phase-annotated action card separately passes its existing 20 checks.
These 135 new regression assertions are not 135 independent empirical
tests of P54. Family choices and finite-array examples remain synthetic.

The recent analyses of Refs.~\cite{PatelShukla,Shukla} motivate checking
finite heavy thresholds before turning a tree Yukawa relation into an
exclusion. Their field content and fitted parameters are not those of P54;
neither their coefficients nor their viability conclusions are imported.
Here the calculable advance is a common phase dictionary and a valid
upper heavy-field matching branch, with explicit algebra for transporting
finite corrections into the light state.

The remaining physical map must contain both-site finite gauge and Yukawa
matching, the appropriate active-field evolution, sequential neutrino and
Weinberg matching, and same-action scalar feedback. An optimization over
an undefined map is neither a global fit nor evidence that the fit domain
is empty. The code therefore distinguishes missing prerequisites from a
computed infeasible point.

\begin{thebibliography}{9}
\bibitem{PatelShukla} K. M. Patel and S. K. Shukla,
\emph{Quantum corrections and the minimal Yukawa sector of SU(5)},
\href{https://arxiv.org/abs/2310.16563}{arXiv:2310.16563v2} (2024),
especially Appendices A and B.
\bibitem{Shukla} S. K. Shukla,
\emph{Revisiting SU(5) Yukawa Sectors Through Quantum Corrections},
\href{https://arxiv.org/abs/2411.06906}{arXiv:2411.06906v2},
revised 25 August 2026.
\bibitem{LWX} M. Luo, H. Wang and Y. Xiao,
\emph{Two-loop Renormalization Group Equations in General Gauge Field Theories},
\href{https://arxiv.org/abs/hep-ph/0211440}{arXiv:hep-ph/0211440}.
\bibitem{MOR} D. Meloni, T. Ohlsson and S. Riad,
\emph{Renormalization Group Running of Fermion Observables in an Extended
Non-Supersymmetric SO(10) Model},
\href{https://arxiv.org/abs/1612.07973}{arXiv:1612.07973}.
\bibitem{Hall} L. J. Hall, \emph{Grand Unification of Effective Gauge Theories},
\href{https://doi.org/10.1016/0550-3213(81)90498-3}{Nucl. Phys. B 178 (1981) 75}.
\end{thebibliography}
\end{document}
